\subsection{The residual chain from an arbitrary Case II point to the seed}
\label{sec:rh-case-II-chain}

We prove the original three-link residual bridge, not yet the
transport of the support of the ordinary Hecke algebra across it.
All group representations over a residue field are continuous
for the discrete topology. Their images are therefore finite.

\begin{lemma}[Auxiliary ramification in a prescribed line]
\label{lem:rh-prescribed-ramification}
Let \(F\) be a number field, \(\ell\) an odd prime, and \(V\)
a finite-dimensional \(\mathbb F_\ell\)-representation of \(G_F\).
Let \(S_0\) contain the primes above \(\ell\) and all the finite
ramified primes of \(V\). Let \(\mathcal Q\) be an infinite set
of other primes at which both \(V\) and \(\mu_\ell\) are trivial
as local Galois modules. Given \(0\ne v_0\in V\), there are
a finite \(S_1\subset\mathcal Q\) and a continuous cocycle
\(c:G_F\to V\), unramified outside \(S_0\cup S_1\), such that
\(c\) is a coboundary at every place of \(S_0\), and, for
some \(q\in S_1\), its restriction to the chosen inertia at \(q\)
is nonzero with values in \(\mathbb F_\ell v_0\).
\end{lemma}
\begin{proof}
Write \(d=\dim_{\mathbb F_\ell}V\). For each finite set \(S_1\)
use the Selmer condition zero at \(S_0\), and at each \(q\in S_1\)
allow all unramified classes and the ramified classes whose
inertia image lies in \(\mathbb F_\ell v_0\). At a prime of
\(\mathcal Q\) the latter condition has dimension \(d+1\):
the trivial module has \(d\) unramified dimensions, and tame
inertia has one independent \(\mathbb F_\ell\)-valued character
because \(Nq\equiv1\bmod\ell\). Thus its local Euler
contribution after subtracting \(h^0\) is one.
At \(S_0\) the contribution is at least \(-d\); at each
infinite place it is at least \(-d\), since \(\ell\) is odd.
The finite Poitou--Tate dimension formula with dual local
annihilators gives the lower bound
\[
 \dim H^1_{\mathcal L}(G_{F,S_0\cup S_1},V)
 \ \ge\ |S_1|-d(1+|S_0|+[F:\mathbb Q]).
\]
Indeed the dual Selmer dimension is nonnegative and the
global \(h^0\) difference is at least \(-d\).
The unramified-outside-\(S_0\) cohomology is finite-dimensional
by the finiteness theorem for Galois cohomology of finite
modules. Take \(|S_1|\) larger than this dimension plus the
displayed constant. If all classes were unramified at every
new prime they would inflate from that fixed finite-dimensional
space, a contradiction. A class with new ramification has the
stated properties; restriction of any of its representatives
to a new inertia group does not depend on the representative
because the action there is trivial.
\end{proof}

\paragraph{Application to a marked frame.}
Fix a tower \(T\) with the original good hypotheses and residue
field \(k\) (it need not be finite). Denote its residual diagonal
characters by \(d_i\) and set \(D_g=\operatorname{diag}(d_i(g))\).
Write \(a_g=d_1(g)/d_0(g)\), \(b_g=d_2(g)/d_0(g)\).
Thus \(d_3/d_0=ab\). These are the characters of
\(-w_{\rm Al},-w_{\rm Ps},-w_{\rm Th}\), respectively \(a,b,ab\).
At the marked element \(s=T.g_0\) all three are different
from one. At any \(g\in T.\mathrm{fr}(W).N\) all \(d_i(g)=1\).

Consider finitely many locally constant \(b\)-cocycles \(y_j\)
killing this set \(N\), with \(y_j(s)=0\). Suppose also that at
each original \(\ell\)-adic place \(y_j(g)=n_{j,v}(1-b_g)\).
We claim that after enlarging \(W\) only by rational primes
satisfying the original \(\texttt{T.Splits}\) condition one can
construct a nonzero normalized \(a\)-cocycle \(x\), and functions
\(z_j\), with the exact identities
\[
\begin{split}
 x(gh)&=x(g)+a_gx(h),\\
 z_j(gh)&=z_j(g)+a_gb_gz_j(h)+y_j(g)a_gx(h),\\
 x(g)&=m_v(1-a_g),\\
 z_j(g)&=o_{j,v}(1-a_gb_g)-y_j(g)a_gm_v
       \qquad(g\in D_v).
\end{split}
\tag{2}
\]
Both \(x,z_j\) vanish on the new killing set and at \(s\).
This is the source's \(\texttt{cupEx}\) assertion in the
needed roots.

For this use the triangular representation on coordinates
\((X,Z_1,\ldots,Z_n)\) given by
\[
 \rho(g)(X,Z_j)_j=(a_gX,\ a_gb_gZ_j+y_j(g)a_gX)_j.
\]
The cocycle law for each \(y_j\) verifies the representation
law. The image of \(\rho\) is finite by continuity and
compactness, even when \(k\) is infinite. Take the finite
\(\mathbb F_\ell\)-submodule spanned by the orbit of
\(v_0=(1,0,\ldots,0)\). It is stable and the first coordinate
of \(v_0\) is nonzero. Set \(S_0\) equal to all places of
\(F\) above the original \(D.S\cup W\).
The action is trivial on inertia outside \(S_0\):
the original good datum has trivial residual characters
on these inertia groups, and these groups are in the original
killing set. Choose an infinite set of rational primes
splitting completely in a finite Galois extension of
\(\mathbb Q\) containing \(F\), the fixed fields of the
kernels of the original residual representation and
\(\omega\), and a normal closure of the field cut out by
\(\rho\). Chebotarev gives infinitely many such primes, none
of which need belong to the original bad set. At all places
over them \(\rho,\mu_\ell\) are trivial. They satisfy the
literal \(\texttt{T.Splits}\) condition.

Apply Lemma~\ref{lem:rh-prescribed-ramification} and enlarge
\(W\) by the rational primes below its \(S_1\). The coordinates
of its cocycle give (2), initially without the marked
normalizations: at an old local group write the local
coboundary as \(\rho(g)c_v-c_v\) and negate the coordinates
of \(c_v\) to obtain \(m_v,o_{j,v}\). On an inertia group
in the new killing set inside \(S_0\) the action is trivial,
so the coboundary vanishes; outside \(S_0\) and the newly
allowed primes the cocycle is unramified. This verifies the
original killing conditions. Some element of new inertia
has \(a_g=1\) and \(x(g)\ne0\).
To normalize at the mark, set
\[
 m_*=\frac{x(s)}{1-a_s},\qquad
 o_{j,*}=\frac{z_j(s)}{1-a_sb_s}.
\]
Subtract \(m_*(1-a_g)\) from \(x(g)\), and subtract
\(o_{j,*}(1-a_gb_g)-y_j(g)a_gm_*\) from \(z_j(g)\).
This is subtraction of the same global coboundary, with
the appropriate sign convention. Replace \(m_v,o_{j,v}\)
by their differences with \(m_*,o_{j,*}\). All identities
persist, and \(y_j(s)=0\) gives the required marked
normalizations. The witness with \(a_g=1\) is unchanged.

\paragraph{Realization as symplectic matrices.}
For a single \(y=y_j\), let
\[
 U(x,y,t)=
 \begin{pmatrix}1&0&0&0\\x&1&0&0\\y&0&1&0\\
                 t&y&-x&1\end{pmatrix}.
\]
This matrix is symplectic for the fixed \(J\) of
Section~\ref{sec:rh-symplectic-induction}, and
\[
 U(x,y,t)U(x',y',t')=
 U(x+x',y+y',t+t'+yx'-xy').
\]
Conjugation by \(D_g\) scales the three coordinates by
\(a_g,b_g,a_gb_g\). The matrices
\[
 r_3(g)=U(x(g),y(g),2z_j(g)-x(g)y(g))D_g,\qquad
 N_{3,v}=U(m_v,n_{j,v},2o_{j,v}-m_vn_{j,v})
\]
therefore give a representation, by (2). On \(D_v\)
they satisfy \(r_3(g)=N_{3,v}D_gN_{3,v}^{-1}\):
the local formula for the third coordinate is
\[
 (1-a_gb_g)(2o_{j,v}-m_vn_{j,v})
       +m_vn_{j,v}(b_g-a_g).
\]
The representation has support \(q_Q\) and the lower
flag has the same support. Its mark is diagonal. The
\(\mathrm{Mu}\) condition holds whenever the original
residual characters are trivial there, which is part of
the good tower conditions (\(\mathrm{Mu}\) lies in inertia
and \(\texttt{Ok.ineL}\) applies). It has exactly the
same normalized \((2,0)\) entry and flag entry as the
input \(y_j,n_{j,v}\), and a nonzero \((1,0)\) entry.
This realizes the source's \(\texttt{gysinAl}\) bridge
simultaneously for all the given \(y_j\), with a single
normalized \(x\) and the same \(m_v\).

\paragraph{Two polynomial links.}
For clarity we record why the requisite interpolations work
over the polynomial ring, including finite residue fields.
A symplectic unipotent matrix of support \(n_P\) has the shape
\[
 \begin{pmatrix}
 1&0&0&0\\x&1&z&0\\y&0&1&0\\t&y&-x+yz&1
 \end{pmatrix}.
\]
In block order \((0,2)\mid(1,3)\) it is block lower triangular;
its diagonal blocks are determined by \(y\). Thus two
residual representations with support \(n_P\) and the
same normalized \((2,0)\) coordinate have identical
diagonal blocks in that order. The same applies to their
local flags of support \(q_Q\) if their \((2,0)\) flag
coordinates agree. In the block triangular algebra the
strict lower-block ideal has square zero. The entrywise
interpolants \(A+X(B-A)\) therefore preserve multiplication
on any pairs with these common diagonal blocks, and preserve
the symplectic identities (the symplectic adjoint is linear,
preserves this block triangular algebra and its square-zero
ideal). The same argument for \(N_v^{-1}r(g)N_v\) shows
that its interpolant is the entrywise interpolation of the
two conjugates, hence retains upper triangularity and the
exact diagonal characters. The mark, the killing conditions
and local \(\mathrm{Mu}\) equations are preserved. The
coefficients are locally constant. The flags still have
support \(q_Q\). This proves \(\texttt{lineP}\) as an
identity in \(k[X]\), not merely on its field-valued points.

For support \(n_M\) use instead block order \((0,1)\mid(2,3)\).
Its diagonal blocks are determined by the diagonal and the
\((1,0)\) coordinate (symplecticity determines \((3,2)\)).
With this coordinate fixed in the representations and flags,
the identical square-zero argument proves \(\texttt{lineM}\).
The second alternative of \(\texttt{Res.shape}\) applies
throughout this link.

\begin{theorem}[The original deep Case II residual chain]
\label{thm:rh-deep-II-chain}
Assume the original \(D.\mathrm{Good}\), Case II, the initial
tower \(T_0\) with \(\texttt{Ok},\texttt{Two},\texttt{TameL}\),
and the original rational pair \(c_0\) with
\(\texttt{LatShape}\) at \(T_0\)'s mark, over the given ambient
residue field. Then there is a tower \(T_1\) with
\(T_0.\mathrm{Le}(T_1)\) and \(\texttt{TameL}(T_1)\) such that
on every deeper \(T\) with \(T_1.\mathrm{Le}(T)\) and
\(\texttt{Deep}(T)\), the integral seed of
Theorem~\ref{thm:rh-seed-II-alt} is linked to every original
residual restriction \(c'_0\) satisfying \(\texttt{OverQ}\)
and \(\texttt{IsRes}\). The three links are polynomial
links of the original frame, and all four endpoints are
\(\texttt{Good}\). In particular this proves the original
\(\texttt{deep\_tower\_residue\_chain\_exists}\) supplier.
\end{theorem}
\begin{proof}
We explain the restriction, auxiliary fields and three
links, retaining the original marked flags.

\emph{Restriction and uniqueness.}
For each chosen local group of an intermediate tower,
choose a Galois element transporting the original place
of the smaller field to that place. Apply its representation
matrix to the original flag and write the resulting flag
in lower-unipotent coordinates. This is possible in the
present support: a matrix with diagonal and support \(n_P\)
factors uniquely as a lower-unipotent matrix of support
\(q_Q\) times an upper triangular matrix. Indeed the
unipotent group displayed above has the normal subgroup
\(z=0\), with quotient the root \((1,2)\); factor off that
root and the invertible diagonal. This factorization also
applies after multiplying by the original lower flag.
Uniqueness follows since a lower unitriangular matrix
that is upper triangular is the identity. A different
transporting element differs by an element of the old
local decomposition group, which stabilizes the original
flag, so the transported flag is independent of the choice.
These are precisely the flags in the original \(\texttt{Over}\)
and \(\texttt{OverQ}\) conditions. The transported graded
characters remain the original diagonal characters because
they are global characters.

Restrict \(c_0\) to \(G_{T_0.F}\) and use these flags.
The representation is normalized at the mark, and the
\(\mathrm{Mu}\) equation follows from the trivial residual
diagonal on inertia at \(\ell\). For the killing set away
from \(\ell\), if the rational prime is outside \(D.T\)
the original rational killing condition applies. At a bad
prime in \(D.S_b\), \(\texttt{Ok.kb}\) makes the original
\(\psi\) trivial on every relative inertia element with
trivial residual diagonal. Transport to the chosen
rational inertia group and use the second part of the
rational killing condition \(I_p\cap\ker\psi\).
This is exactly the passage from the rational frame
(\(\ell\) and \(D.S_b\) excluded in its first killing set)
to the tower frame. The same argument works on every
further tower. Thus this restriction is the unique
\(\texttt{OverQ}\) residual restriction, denoted \(p_0\)
on \(T_0\). Its normalized coordinate
\(y(g)=p_0.r(g)_{20}/d_0(g)\) is a \(b\)-cocycle by
the displayed \(n_P\) matrix law. Its local value is
\(p_0.n(v)_{20}(1-b_g)\), by comparing the first columns
of \(r(g)N_v=N_v(N_v^{-1}r(g)N_v)\). It kills \(N\)
and is zero at the mark.

Two elementary facts about restriction will be used.
First, a cocycle \(f\) for a global residual character
\(\chi\), nontrivial at a mark in a normal subgroup \(H\),
cannot vanish identically on \(H\) unless it vanishes
identically on the whole group: for \(h\in H\),
\(f(hg)=\chi(h)f(g)\) and
\(f(g(g^{-1}hg))=f(g)\).
Apply this to the \((2,0)\) and \((1,2)\) normalized
coordinates in \(n_P\), which are cocycles. When the
\((2,0)\) coordinate vanishes identically, the \((1,0)\)
coordinate is a cocycle as well.
Second, when a flag is transported by \(r(g)\), its
new \((2,0)\) lower coordinate is
\[
 \frac{r(g)_{20}+r(g)_{22}N_{20}}{r(g)_{00}}.
 \tag{3}
\]
This follows from the first column in the triangular
factorization, since row two has only its diagonal and
column-zero entries in support \(n_P\). Thus equality
of the global and local \((2,0)\) coordinates is
preserved by simultaneous restriction; if the global
coordinate is zero and the local one nonzero, the latter
stays nonzero.

These facts show directly that the initial
\(\texttt{LatShape}\) alternatives imply the following
\(\texttt{Good}\) alternatives on any tower with its
marked element. If \(\operatorname{Ent}(-w_{\rm Ps})\)
holds initially, both it and \(\operatorname{Ent}(w_{\rm Be})\)
survive, by the cocycle fact. Otherwise the support is
\(v_V\) and the \((2,0)\) flag coordinate is nonzero
by (3). The \((1,0)\) coordinate cannot be identically
zero initially: if so, use the local first-column identity
with \(\chi_0=\chi_2\) and \(N_{20}\ne0\); it would give
\[
 r(g)_{12}=(N_{10}/N_{20})(d_2(g)-d_1(g))
\]
on the original rational decomposition group, contrary
to the last non-coboundary condition in \(\texttt{LatShape}\).
This \((1,0)\) coordinate is then a nonzero cocycle and
survives restriction by the first fact.
Both initial alternatives are therefore good.

\emph{The auxiliary tower.}
Apply the preceding cocycle construction on \(T_0\), with
\(y\) from \(p_0\) (starting \(W=\varnothing\)).
It gives the bridge \(r_3,N_{3,v}\) of support \(q_Q\)
with the same \((2,0)\) coordinates and a nonzero
\((1,0)\) coordinate, over an enlarged set \(W\)
of split rational primes. In particular the auxiliary
primes are chosen before the deeper tower or its seed.

For each \(w\in W\) one has \(w\ne\ell\) and
\(w\equiv1\bmod\ell\). Let \(F_1\) be the compositum
of \(T_0.F\) with the unique degree-\(\ell\) subfield
of each \(\mathbb Q(\zeta_w)\). These fields are totally
real because their degree is odd, and unramified above
\(\ell\). The compositum is abelian over \(\mathbb Q\);
over \(T_0.F\) its Galois group is an elementary abelian
\(\ell\)-group. Keep the old \(K\), ordering, and choose
the mark \(s^\ell\). Choose the specified local places
above \(\ell\) with the original restriction to \(K\)
(which is split there). These give the original
\(\texttt{Le}\) condition. The tame condition at \(\ell\)
persists since all added fields are unramified there.

For any \(\texttt{Le}\) extension, the inclusion of
absolute Galois groups is normal of \(\ell\)-power index
and the new mark is an \(\ell\)-power of the old one.
All the original \(\texttt{Ok}\) and \(\texttt{Two}\)
conditions persist. Here are the details concerning
the existential clauses: replace a witness by the
power \(\ell^j\) supplied by \(\texttt{Le}\), which
preserves every inequality between finite residual
character values, since their orders are prime to
\(\ell\); this applies also inside the chosen inertia
or decomposition groups. The ordering and the
restriction conditions persist. For \(\texttt{Ok.g3'}\),
a quadratic permutation character on the smaller
CM Galois subgroup equals a ratio of two of the
original characters (fix one index). Extend it by
this ratio to the old CM Galois group. Taking
\(\ell^j\)-powers proves that the extension is quadratic,
trivial on the old \(\ker\omega\), and permutes the
same characters. The old condition forces it trivial.
All other clauses are restrictions of universal
conditions or fixed field and ordering conditions;
the real degree stays even. This also proves the
needed composition of tower inclusions. In particular
each deeper \(T\) in the theorem satisfies
\(\texttt{Ok},\texttt{Two},\texttt{Deep}\), so the seed
theorem applies.

On inertia at any new \(w\), the diagonal characters
are trivial. The bridge takes values in \(U(x,y,t)\),
which has exponent \(\ell\): its \(n\)-th power is
\(U(nx,ny,nt)\). Its image is a finite \(\ell\)-group
(the group of such matrices over the finite span of
the coordinate values is an \(\ell\)-group, or use
the central extension by additive groups).
Wild inertia at \(w\) has trivial image since it is
pro-\(w\). Tame inertia has procyclic pro-\(\ell\)
quotient, and its unique quotient of order \(\ell\)
is killed by the degree-\(\ell\) subfield of
\(\mathbb Q(\zeta_w)\), which is totally ramified
at \(w\). Since \(w\) splits in \(T_0.F\), this also
kills its restriction over \(T_0.F\). Hence the bridge
restricted to \(G_{F_1}\), and to every deeper field,
satisfies the killing condition with \(W=\varnothing\).
This argument works at \(\ell=3\) because the bridge
has nilpotence class two, not three.

Transport the flags of the bridge and \(p_0\) to the
deeper \(T\) as above. The bridge flag is transported
inside \(q_Q\), since its global support is \(q_Q\).
Equation (3) shows that the two residual restrictions
have the same \((2,0)\) coordinates globally and at
the flags. The nonzero \(a\)-cocycle \(x\) survives by
normal restriction and the nontriviality of \(a\) at
the new mark. Denote the resulting pairs by \(p'_0,c_3\);
by uniqueness \(p'_0\) is every admissible \(c'_0\).
If the \((2,0)\) global coordinate of \(c_3\) is nonzero,
the first good alternative holds with \(-w_{\rm Al}\).
If it is zero, that coordinate of \(p'_0\) is zero
and all the \((2,0)\) flag coordinates are nonzero,
by the initial good alternatives; \(c_3\) has support
\(v_V\) and is again good.

\emph{The three links.}
Project \(c_3.r\) and its flags onto their diagonal
blocks in order \((0,1)\mid(2,3)\). This is a group
homomorphism on the lower block parabolic and preserves
the symplectic multiplier. Since the original local
conjugates are both lower and upper triangular they
are diagonal, so the projected local flags still give
the same diagonal. The resulting representation has
only the root \(-w_{\rm Al}\) off the diagonal and
retains its nonzero \((1,0)\) entry. Call the projected
flags \(N_v^0\). Write \(Q\) for the new integral
Case II seed reduced into the given residue field, and
\(q_v=Q.n(v)_{20}\ne0\). Replace each projected flag by
\[
 N_v^0(1+q_v(E_{20}+E_{31})).
\]
The right factor commutes with the residual diagonal
on the local group: \(\texttt{Two.atL}\) gives
\(d_0=d_2,\ d_1=d_3\). Hence the flag equations
remain valid. This matrix is lower symplectic of
support \(q_Q\), with the original projected
\((1,0)\) coordinate and \((2,0)=q_v\).
Call this pair \(c_4\). It is good by the
\(\texttt{InV}\) and \(-w_{\rm Al}\) alternative.
The seed itself is good: the entry alternative in
Theorem~\ref{thm:rh-seed-II-alt} either gives
\(-w_{\rm Al}\), or both \(w_{\rm Be},-w_{\rm Th}\);
in the latter case \(\texttt{Two.bigL}\) supplies
the required local nontriviality of \(w_{\rm Al}\).

Finally use the polynomial links already proved:
\[
 p'_0\ \xrightarrow{\ \texttt{lineP}\ }\ c_3\
 \xrightarrow{\ \texttt{lineM}\ }\ c_4\
 \xrightarrow{\ \texttt{lineP}\ }\ \bar Q .
\]
The first has equal global and flag \((2,0)\) entries
by construction and transport. The second has equal
global and flag \((1,0)\) entries by the projection.
The third has zero global \((2,0)\) entries on both
sides and equal flag entries \(q_v\).
All are in the original frame and all four endpoints
are good, as verified above. This proves its exact
chain condition.
\end{proof}

The theorem constructs the residual links for every deeper
tower using auxiliary primes chosen at the initial tower.
It does not say that the three links lie in the support of
the rank-four ordinary automorphic quotient. That further
implication is a characteristic-prime component-propagation
statement and is kept separate.
